\section{A uniform arithmetic lower bound from the cosine pair test}
\label{sec:cosine-lower-rate}

The preceding arithmetic rank construction can also be used in a
uniform sampling estimate. We retain the exact constant
\[
 d=d_{\rm MT}=\frac54-\frac1{2\sqrt2}\cot\frac1{\sqrt2},
 \qquad p_{\rm MT}=d^{-1}=1.1958141\ldots.
\]
Mori proved the lower rate with exponent \(6/5\) by a uniform
interval version of Montgomery's pair argument and quantitative
sampling \cite[Theorem~4.3, pp.~11--13]{mori-rh-alone2026}.
We give the cosine-weight variant, including the uniform errors and
the precise sampling reduction. The zero-location input throughout
is the established theorem already used for the Weil sampling identity.

\begin{theorem}[An all-window lower rate]
\label{thm:cosine-lower-rate}
There are constants \(C,L_0>0\) such that, for \(L\ge L_0\),
\[
 a_1(L)\ge
 \exp\{-C e^{2p_{\rm MT}L}L(1+\log L)\}.
\]
The same lower bound holds for \(b_L\) and \(a_2(L)\).
Equivalently, with \(\mu=e^{2L}\), the exponent is
\(C\mu^{p_{\rm MT}}\log\mu(1+\log\log\mu)\).
\end{theorem}

\subsection{The uniform pair estimate}

Put \(c=1/\sqrt2\), and let
\[
 f_0(v)=\frac{\cos(\sqrt2v)}{\sqrt2\sin c}
              \mathbf1_{[-1/2,1/2]}(v),\qquad
 w=f_0*\widetilde f_0,
\]
where \(\widetilde f_0(v)=f_0(-v)\). Write
\[
 A=\int f_0^2=w(0),\qquad
 B=\int |v|w(v)\,dv.
\]
Then \(w\) is even and nonnegative, is supported in \([-1,1]\),
\(\int w=1\), and \(0\le w\le A\) by Cauchy's inequality.
The exact energy calculation in Section~\ref{sec:survival} gives
\begin{equation}\label{eq:cosine-energy}
 A+B=\frac12+c\cot c=3-2d.
\end{equation}
We shall use the elementary bounds
\begin{equation}\label{eq:cosine-coarse}
 1\le A\le\frac{51}{50},\qquad
 0<B\le\frac13,\qquad
 \frac56<d\le\frac78.
\end{equation}
Here \(A\ge1\) follows from \(\int f_0=1\). The variance of
\(f_0\) under uniform measure on its support is \(A-1\), and is
at most one quarter of its squared range. That range is
\(\tan(c/2)/\sqrt2\). Since
\(\tan(c/2)\le (c/2)/(1-c^2/8)=8/(15\sqrt2)\),
we get \(A-1\le4/225<1/50\).
The Taylor bounds
\(\cos c\le73/96\), \(\sin c\ge11c/12\) show
\(c\cot c\le73/88<5/6\). Thus \(A+B<4/3\),
which gives \(B<1/3\) and \(d>5/6\).
Conversely \(c\cot c\ge\cos c\ge3/4\), giving \(d\le7/8\).

Let \(J\) be any bounded interval, and define
\[
 \ell=|J|,\qquad H=\max\{2,\sup_{t\in J}|t|\},
 \qquad L_J=\log(2+H).
\]
Let \(N_J\) count all zero labels in \(J\), including multiplicity,
and let \(P_h(J)\) count unordered pairs of distinct labels at distance
less than \(h\); equal ordinates are allowed. There is an absolute
constant \(C_0\) such that, whenever
\(\kappa\) is sufficiently large,
\(L_J\ge\kappa\), \(\ell\ge e^\kappa\), and
\(h=\kappa^{-3/2}\),
\begin{equation}\label{eq:cosine-uniform-pair}
 N_J+2P_h(J)
 \le\frac{\ell}{2\pi}
       \left(\frac{A L_J^2}{\kappa}+B\kappa\right)
          (1+C_0/\kappa).
\end{equation}
All constants are independent of the position of \(J\).

Here are the arithmetic inputs and the proof of this uniformity.
For \(x\ge1\) set
\[
 \mathcal L_x(t)=2\sum_\gamma
           \frac{m_\gamma x^{i\gamma}}{1+(t-\gamma)^2}.
\]
Montgomery's explicit-formula decomposition is
\(\mathcal L_{e^u}=\mathcal A_u+\mathcal B_u+\mathcal E_u\),
where
\[
 \mathcal A_u(t)=-e^{iut}\sum_{n\ge1}
       \frac{\Lambda(n)}{n^{1/2+it}}
                  \min\{ne^{-u},e^u/n\},
 \qquad
 \mathcal B_u(t)=e^{-u+iut}
           \left(-\frac{\zeta'}\zeta(-1/2+it)\right).
\]
For an interval of length \(r\) with
\(\log(|t|+2)\le V\), the following bounds hold with absolute
constants:
\begin{equation}\label{eq:cosine-Montgomery-input}
 \begin{split}
 \int|\mathcal A_u(t)|^2dt
    &\le r(u+C)+Ce^u(u+1),\\
 |\mathcal B_u(t)|&\le e^{-u}(V+C),\\
 \int_{\mathbb R}|\mathcal E_u(t)|^2dt&\le Ce^u.
 \end{split}
\end{equation}
These precise interfaces, with the decomposition derived from
Landau's formula, are proved in
\cite[Lemmas~3.3--3.5, pp.~6--7]{mori-rh-alone2026}.
For clarity, the mean-value step uses the Montgomery--Vaughan
weighted Hilbert inequality on frequencies \(\log n\), whose local
spacing is at least \(1/(2n)\). The coefficient square sums satisfy
\(\sum|a_n(e^u)|^2\le u+C\) and
\(\sum n|a_n(e^u)|^2\le Ce^u(u+1)\), by partial summation
from the prime number theorem. The functional equation gives the
second bound, and the pole and trivial-zero terms give the third.
No close-pair hypothesis is used in these bounds.

Enlarge \(J\) by \(Y=\ell^{1/4}\) on each side, obtaining
\(J^+\). Put \(\ell^+=\ell+2Y\) and
\(L^+=\log(2+H+Y)\). On functions of \((u,t)\) define the
squared Hilbert norm
\[
 \|F\|_w^2=\frac2\kappa\int_0^\kappa w(u/\kappa)
                       \int_{J^+}|F(u,t)|^2dt\,du.
\]
The prime term in \eqref{eq:cosine-Montgomery-input} gives
\[
 \|\mathcal A\|_w^2
       \le \ell^+(B\kappa+C)+Ce^\kappa.
\]
Indeed \((2/\kappa)\int_0^\kappa u w(u/\kappa)du=B\kappa\),
and
\((2/\kappa)\int_0^\kappa w(u/\kappa)e^u(u+1)du
 \le 2Ae^\kappa\).
Similarly,
\[
 \|\mathcal B\|_w^2
       \le\frac{A\ell^+(L^++C)^2}{\kappa},
 \qquad
 \|\mathcal E\|_w^2\le Ce^\kappa/\kappa.
\]
It is essential to integrate the prime--archimedean cross term
before applying an inequality to the full \(u\)-norm. Cauchy's
inequality in \(t\) alone bounds its absolute value by
\[
 \frac{C(L^++C)\sqrt{\ell^+}}{\kappa}
 \int_0^\kappa e^{-u}
   \left(\sqrt{\ell^+(u+C)}+e^{u/2}\sqrt{u+1}\right)du
 \le\frac{C\ell^+(L^++C)}{\kappa}.
\]
The integrals on \((0,\infty)\) are finite absolute constants.

Let
\(Q=\ell(A L_J^2/\kappa+B\kappa)\).
Since \(H\ge\ell/2\),
\(L^+-L_J\le2\ell^{-3/4}\) and
\(\ell^+/\ell=1+2\ell^{-3/4}\).
Using \(L_J\ge\kappa\) and \(\ell\ge e^\kappa\), the preceding
bounds yield
\[
 \|\mathcal A+\mathcal B\|_w^2
       \le Q(1+C/\kappa),\qquad
 \|\mathcal E\|_w^2\le CQ/\kappa^2.
\]
The triangle inequality therefore gives
\begin{equation}\label{eq:cosine-average-upper}
 \|\mathcal L\|_w^2\le Q(1+C/\kappa).
\end{equation}

Expanding the square gives the exact positive pair average
\[
 \|\mathcal L\|_w^2
 =4\sum_{\gamma,\gamma'}m_\gamma m_{\gamma'}
 \widehat w_1\bigl(\kappa(\gamma-\gamma')\bigr)
 \int_{J^+}\frac{dt}
 {(1+(t-\gamma)^2)(1+(t-\gamma')^2)},
\]
where \(\widehat w_1(v)=\int w(y)e^{ivy}dy
=|\int f_0(y)e^{ivy}dy|^2\ge0\).
The double sum converges absolutely on the bounded integration
interval. For \(|v|\le\kappa h=\kappa^{-1/2}\),
\(\widehat w_1(v)\ge1-v^2/2\), since \(w\) is a probability
density supported in \([-1,1]\). For \(\gamma,\gamma'\in J\),
the Cauchy-kernel integral is at least
\[
 \frac{2\pi}{4+(\gamma-\gamma')^2}-\frac{2}{3Y^3}.
\]
Thus every retained ordered pair at distance less than \(h\),
including equal labels, contributes at least
\(2\pi(1-C/\kappa)\). Discarding all other nonnegative terms
and using \eqref{eq:cosine-average-upper} proves
\eqref{eq:cosine-uniform-pair}.

\subsection{Packing with the cosine constant}

The uniform interval count from Riemann--von Mangoldt is
\[
 N_J\ge\ell\left(\frac{L_J}{2\pi}-b_0\right)-c_0 L_J
\]
for absolute constants \(b_0,c_0\). Its proof averages \(\log|t|\)
on \(J\), with endpoint errors \(O(L_J)\); this gives a constant
loss in the average logarithm even for intervals containing zero
\cite[Lemma~2.2, pp.~7--8]{mori-sharp2026}.
Let \(N_0=\ell L_J/(2\pi)\), \(\beta=\kappa/L_J\), and
\(\rho=A/\beta+B\beta\). Then
\[
 N_J\ge N_0(1-C/\kappa),\qquad
 N_J+2P_h(J)\le\rho N_0(1+C/\kappa).
\]
In the close-pair graph, deletion of one endpoint per edge gives
an independent set of size at least \(N_J-P_h(J)\). Ordering its
vertices uniformly at random and retaining each vertex earlier than
all its neighbors gives expected size
\(\sum_v(1+\deg v)^{-1}\ge N_J^2/(N_J+2P_h(J))\).
Therefore the maximal size \(p_h(J)\) of an \(h\)-separated
subset is bounded below by the larger of these two expressions.

The exact scalar comparison needed here is
\begin{equation}\label{eq:cosine-scalar}
 \max\{3/2-\rho/2,1/\rho\}\ge d\beta
                     \qquad(0<\beta\le1).
\end{equation}
For \(\beta\ge3/5\), the difference in its first branch is
\[
 \frac{(1-\beta)((3-A)\beta-A)}{2\beta}\ge0,
\]
by \(A\le51/50\). For \(\beta\le3/5\),
\[
 \rho^{-1}=\frac\beta{A+B\beta^2}
       \ge\frac{50}{57}\beta>\frac78\beta\ge d\beta.
\]
In the first range \(\rho\) is bounded above and \(\beta\)
is bounded away from zero, so the preceding relative errors cost
only \(C/\kappa\) relative to \(d\beta\). In the second range
the reciprocal branch carries a relative \(O(1/\kappa)\) error
directly. We obtain, uniformly in the position of \(J\),
\begin{equation}\label{eq:cosine-separated}
 p_h(J)\ge\frac{d\kappa}{2\pi}\ell(1-C_1/\kappa),
 \quad L_J\ge\kappa,\quad\ell\ge e^\kappa,
 \quad h=\kappa^{-3/2}.
\end{equation}

\subsection{Quantitative sampling and the exact window map}

We recall the quantitative sampling consequence in the form needed
here, and specify its classical input. If a \(q\)-separated set
\(\Lambda\subset\mathbb R\) satisfies
\[
 \#(\Lambda\cap I)\ge(1+\epsilon)|I|-D
       \quad\hbox{for every bounded interval }I,
 \qquad0<q,\epsilon\le1,
\]
then its lower sampling constant on \(PW_\pi\) satisfies
\begin{equation}\label{eq:cosine-VH}
 \inf_{0\ne F\in PW_\pi}
 \frac{\sum_{x\in\Lambda}|F(x)|^2}{\|F\|_2^2}
 \ge\exp\{-C(D+1)(1+\log(1/q)+\log(1/\epsilon))\}.
\end{equation}
Here is the reduction to Voss's theorem. Put
\(r=(D+1)/(2(1+\epsilon))\) and \(\rho_0=1+\epsilon/2\).
For any finite set of integers, the union of the intervals
\([n/\rho_0-r,n/\rho_0+r]\) contains enough points of
\(\Lambda\) for an injective matching: a component containing
\(m\) centers has length at least \((m-1)/\rho_0+2r\), and the
count lower bound is at least \(m\). Finite Hall matching and a
locally finite diagonal limit give points \(x_n\in\Lambda\)
with \(|x_n-n/\rho_0|\le r\).

Set \(t_n=\rho_0x_n\),
\(M=\lfloor\rho_0r\rfloor+1\), and
\(\delta_0=\tfrac12\min(1,\rho_0q)\).
Then \(|t_n-n|<M\) and the separation is strictly greater than
\(\delta_0\). For \(F\in PW_\pi\),
\(G(x)=F(x/\rho_0)\) has type \(\tau=\pi/\rho_0<\pi\)
and \(\|G\|_2^2=\rho_0\|F\|_2^2\).
The sampling sequence is square summable: integrating the fundamental
theorem of calculus over the disjoint intervals of length
\(\delta_0\) centered at \(t_n\) gives
\[
 \sum_n|G(t_n)|^2\le
 \frac{2}{\delta_0}\|G\|_2^2+\delta_0\|G'\|_2^2
 \le(2/\delta_0+\delta_0\tau^2)\|G\|_2^2.
\]
Voss's published quantitative bound
\cite[Theorem~2 and Corollary~1, pp.~195--196]{voss1999} is
\[
 \sum_n|G(t_n)|^2\ge\frac{H_V^2}{2M+1}\|G\|_2^2,
\]
where
\[
 H_V=\delta_0^{2M}
 \left(\frac1{\pi-\tau}+\frac{21}{10}\right)^{-4M}
 \left(\frac2\pi\right)^{2M-1}
 \frac{M}{3e^{4M}}\left(4+\frac1M\right)^{-4M}.
\]
Since \(M=O(D+1)\), \(\delta_0\ge q/2\), and
\(\pi-\tau\ge\pi\epsilon/3\), taking logarithms proves
\eqref{eq:cosine-VH}. This reproduces the Voss--Hall reduction of
\cite[Proposition~2.4, p.~8]{mori-sharp2026}, with strict endpoint
parameters and with only an absolute constant retained.

There is one \(h\)-separated set of zero ordinates \(\Gamma_h\)
whose count in every bounded interval is at least \(p_h(J)-1\).
Indeed, left-to-right greedy selection on a finite interval is
optimal there. Its restriction to a subinterval loses at most one
point compared with starting the greedy selection at that
subinterval. Letting the finite endpoints tend to infinity along a
locally finite diagonal subsequence gives the asserted set.

Now put \(s=2L=\log\mu\), and choose
\[
 \kappa=\frac{s+1+2\pi}{d}+C_1+1,
 \qquad h=\kappa^{-3/2}.
\]
In particular \(d(\kappa-C_1)=s+1+2\pi+d\).
For sufficiently large \(s\), equation~\eqref{eq:cosine-separated}
gives
\(p_h(J)\ge((s+1)/(2\pi)+1)|J|\) whenever
\(L_J\ge\kappa\) and \(|J|\ge e^\kappa\).
Write \(r_s=(s+1)/(2\pi)\) and \(D=2r_s e^\kappa\).
On these long intervals the one-point greedy loss is paid by the
extra density one. On shorter intervals with \(L_J\ge\kappa\),
the possible deficit is at most \(r_s e^\kappa\). Intervals
with \(L_J<\kappa\) lie inside \((-e^\kappa,e^\kappa)\),
so their possible deficit is at most \(D\). Consequently
\[
 \#(\Gamma_h\cap J)\ge r_s|J|-D
                      \quad\hbox{for every bounded interval }J.
\]
Scale to \(\Lambda=s\Gamma_h/(2\pi)\). This set has separation
\(q=sh/(2\pi)\), satisfies the sampling count with
\(\epsilon=1/s\), and has
\[
 D=O(s\mu^{1/d}),\qquad
 1+\log(1/q)+\log(1/\epsilon)=O(1+\log s).
\]
Equation~\eqref{eq:cosine-VH} bounds its sampling constant below by
\(\exp\{-Cs\mu^{1/d}(1+\log s)\}\).

For the final normalization, take a compact smooth physical function
\(f\) supported in \((-L,L)\), and set
\(F(x)=\widehat f(\pi x/L)\). Then \(F\in PW_\pi\), and
Plancherel gives \(\|F\|_2^2=2L\|f\|_2^2=s\|f\|_2^2\).
The Weil sampling formula, with positive multiplicities,
gives exactly
\[
 q(f)\ge\sum_{\gamma\in\Gamma_h}|\widehat f(\gamma)|^2
       =\sum_{x\in\Lambda}|F(x)|^2.
\]
The compact smooth core and the variational characterization now
prove Theorem~\ref{thm:cosine-lower-rate} for \(a_1(L)\), after
absorbing the harmless factor \(s\) into the constant.
Since \(a_1(L)\le b_L\le a_2(L)\), its other assertions follow.

This lower rate concerns the Weil form at every
sufficiently large window. Its exponent satisfies
\(1<p_{\rm MT}<6/5\). The required denominator estimate in
Section~\ref{sec:theta-coercivity} has exponential
\(e^{-4\pi e^{2L}}\) and a positive polynomial gain. The bound
proved here does not supply that estimate: inserting it into the
prolate trial bound leaves the positive term
\(C e^{2p_{\rm MT}L}L(1+\log L)\) in the exponent of the resulting
upper ratio bound. No vanishing ratio is asserted from this rate.
